\chapter{自由幺半群,自由群}

\section{自由原群}

记号约定:$X,Y,Z$是集合,$N$是原群.

\begin{definition}{阶次分量}{graded-components}
	递归地定义集合序列$\left(M_{n}\left(X\right)\right)_{n=1}^{\infty}$如下:
	\begin{align}
		M_{n}\left(X\right)=
		\begin{cases}
			X,&n=1\\
			\coprod\limits_{p=1}^{n-1}\left(M_{p}\left(X\right)\times M_{n-p}\left(X\right)\right),&n\geq 2
		\end{cases}
	\end{align}
	对诸$p,q\in\N_{>0}$,记$\phi:M_{p}\left(X\right)\times M_{q}\left(X\right)\hookrightarrow M_{p+q}\left(X\right)$为典范投射.
\end{definition}

\begin{definition}{$\Mag\left(X\right)$的构造}{construction-of-mag-x}
	$\left(M_{n}\left(X\right)\right)_{n=1}^{\infty}$如定义(\ref{def:graded-components})所述.定义$\Mag\left(X\right)\coloneq\coprod\limits_{n=1}^{\infty}M_{n}\left(X\right)$.
\end{definition}

\begin{proposition}{长度的良定义性}{}
	对诸$\omega\in\Mag\left(X\right)$,存在唯一的$n\in\N_{>0}$使得$\omega\in M_{n}\left(X\right)$.
\end{proposition}

\begin{definition}{长度函数,长度}{}
	定义$\ell:\Mag\left(X\right)\to\N_{>0}$如下:对诸$\omega\in\Mag\left(X\right)$,$\omega\in M_{\ell\left(\omega\right)}\left(X\right)$.我们称$\ell$为长度函数,对诸$\omega\in X$,称$\ell\left(\omega\right)$为其长度.
\end{definition}

\begin{definition}{$\Mag\left(X\right)$上的运算}{}
	定义$\Mag\left(X\right)$上的运算$\cdot$为$\left(\omega,\omega^{\prime}\right)\mapsto\phi_{\ell\left(\omega\right),\ell\left(\omega^{\prime}\right)}\left(\omega,\omega^{\prime}\right)$.
\end{definition}

\begin{proposition}{$\ell$的加性,$M\left(X\right)$的唯一分解性}{}
	\begin{enumerate}
		\item $\ell$是半群同态.
		\item $\forall \omega\in\Mag\left(X\right)\left(\ell\left(\omega\right)\geq 2\implies\exists!\omega^{\prime},\omega^{\prime\prime}\in \Mag\left(X\right)\left(\omega=\omega^{\prime}\cdot\omega^{\prime\prime}\right)\right)$.
	\end{enumerate}
\end{proposition}

\begin{definition}{自由原群}{}
	配备运算$\cdot$的$\Mag\left(X\right)$称为自由原群.
\end{definition}

\begin{proposition}{自由原群的泛性质}{}
	设$M$是原群,$f\in\Hom_{\mathsf{Set}}\left(X,M\right)$,则存在唯一的$\phi\in\Hom_{\mathsf{Mag}}\left(\Mag\left(X\right),M\right)$使下图交换(其中,$\iota$是典范嵌入).
	\begin{center}
	\begin{tikzcd}
		X \arrow[rr, "\iota", hook] \arrow[rd, "f"'] &   & \Mag\left(X\right) \arrow[ld, "\exists!\phi", dashed] \\
		& M &                                                         
	\end{tikzcd}
	\end{center}
\end{proposition}

\begin{definition}{诱导映射$\Mag\left(u\right)$}{}
	设$u\in\Hom_{\mathsf{Set}}\left(X,Y\right)$.定义$\Mag\left(u\right)\in\Hom_{\mathsf{Mag}}\left(\Mag\left(X\right),\Mag\left(Y\right)\right)$为使下图交换的映射.
	\begin{center}
		\begin{tikzcd}
			X \arrow[rr, "\iota", hook] \arrow[d, "u"'] &  & \Mag\left(X\right) \arrow[d, "\exists!\Mag\left(u\right)", dashed] \\
			Y \arrow[rr, "\iota^{\prime}"', hook]       &  & \Mag\left(Y\right)                                                
		\end{tikzcd}
	\end{center}
\end{definition}

\begin{proposition}{$\Mag\left(-\right)$的函子性}{}
	设$u\in\Hom_{\mathsf{Set}}\left(X,Y\right),v\in\Hom_{\mathsf{Set}}\left(Y,Z\right)$,则下图交换.
	\begin{center}
		\begin{tikzcd}
			X & {\Mag\left(X\right)} \\
			Y & {\Mag\left(Y\right)} \\
			Z & {\Mag\left(Z\right)}
			\arrow["\iota", hook, from=1-1, to=1-2]
			\arrow["u", from=1-1, to=2-1]
			\arrow["{v\circ u}"'{pos=0.5}, bend right=40, from=1-1, to=3-1]
			\arrow["{\Mag\left(u\right)}"', from=1-2, to=2-2]
			\arrow["{\Mag\left(v\circ u\right)}"{pos=0.5}, shift left=2pt, bend left=50, from=1-2, to=3-2]
			\arrow["{\iota^{\prime}}", hook, from=2-1, to=2-2]
			\arrow["v", from=2-1, to=3-1]
			\arrow["{\Mag\left(v\right)}"', from=2-2, to=3-2]
			\arrow["{\iota^{\prime\prime}}", hook, from=3-1, to=3-2]
		\end{tikzcd}
	\end{center}
\end{proposition}

\begin{proposition}{$\Mag\left(-\right)$保持单性和满性}{}
	设$u\in\Hom_{\mathsf{Set}}\left(X,Y\right)$.若$u$是单射(相对地,满射,双射),则$\Mag\left(u\right)$是单射(相对地,满射,双射).
\end{proposition}

\begin{proposition}{子原群的嵌入}{}
	设$S\subseteq X$,$i:S\hookrightarrow X$是典范嵌入,则下图交换.因此$\Mag\left(S\right)$在同构意义下是$S$在$\Mag\left(X\right)$中生成的子原群.
	\begin{center}
		\begin{tikzcd}
			S & {\Mag\left(X\right)} \\
			X & {\Mag\left(Y\right)}
			\arrow["\iota", hook, from=1-1, to=1-2]
			\arrow["i"', hook, from=1-1, to=2-1]
			\arrow["{\Mag\left(i\right)}", hook', from=1-2, to=2-2]
			\arrow["{\iota^{\prime}}", hook, from=2-1, to=2-2]
		\end{tikzcd}
	\end{center}
\end{proposition}

本节接下来的内容中,设$\left(\left(u_{\alpha},v_{\alpha}\right)\right)_{\alpha\in I}$是$\Mag\left(X\right)$的元素族,$\sim$是其生成的相容等价关系,$\left(n_{x}\right)_{x\in X}$是$N$的元素族.

\begin{proposition}{}{}
	设$\pi:\Mag\left(X\right)\twoheadrightarrow\Mag\left(X\right)/\sim$是典范满射,则$\Mag\left(X\right)/\sim$作为原群由$\pi\left(X\right)$生成.
\end{proposition}

\begin{proposition}{唯一扩张定理}{}
	设$k\in\Hom_{\mathsf{Mag}}\left(\Mag\left(X\right),N\right)$由映射$X\to N,x\mapsto n_{x}$延拓得到.若$\forall\alpha\in I\left(k\left(u_{\alpha}\right)=k\left(v_{\alpha}\right)\right)$,则下图交换.
	\begin{center}
		\begin{tikzcd}
			{\Mag\left(X\right)} && {\Mag\left(X\right)/\sim} \\
			& N
			\arrow["\pi", two heads, from=1-1, to=1-3]
			\arrow["k"', from=1-1, to=2-2]
			\arrow["{\exists!f}", dashed, from=1-3, to=2-2]
		\end{tikzcd}
	\end{center}
\end{proposition}












